\chapter{The 2D Virtual Office}
\label{ch:skyoffice}

The 2D office is what makes \voname{} a \emph{place} rather than a set of tools.
This chapter describes it in full: its client/server split, the shared type
contract that keeps the two in step, the multiplayer synchronisation model that
keeps many avatars moving smoothly, the interactive objects that populate the
world, and the peer-to-peer stack that carries proximity voice, video, and screen
sharing. The component is an adapted fork of the open-source \emph{SkyOffice}
project, extended and integrated as one of \voname{}'s three clients.

\section{Overview and Structure}

The 2D office is a small but complete \emph{multiplayer game}. It is organised as
a monorepo of three workspaces (Table~\ref{tab:sky-structure}): a browser
\emph{client} that renders the world and the user interface, an authoritative
\emph{server} that holds the shared state, and a \emph{types} package of shared
definitions imported by both. Its one-line technology summary is: Phaser for the
game, React and Redux for the interface, Colyseus for multiplayer state, and
PeerJS/WebRTC for media.

\begin{table}[htbp]
\centering
\caption{Structure of the 2D office component.}
\label{tab:sky-structure}
\small
\begin{tabularx}{\textwidth}{L{2.0cm} X L{3.6cm}}
\toprule
\textbf{Workspace} & \textbf{Role} & \textbf{Stack} \\
\midrule
\code{client/} & Renders the world and UI; captures input; runs media & TypeScript, Phaser~3, React~18, Redux Toolkit, PeerJS, Vite \\
\code{server/} & Authoritative multiplayer state and message handling & TypeScript, Colyseus, Express, bcrypt \\
\code{types/} & Shared enums, message codes, and state interfaces & TypeScript \\
\bottomrule
\end{tabularx}
\end{table}

\section{The Shared Type Contract}

Because client and server are separate programs that must agree exactly on the
shape of the world and the messages they exchange, the \code{types/} package is
the linchpin: both sides import it, so a change to a message code or a state field
is enforced by the compiler across the whole system. Its central definitions are:

\begin{itemize}
  \item \code{IOfficeState} --- the synchronised room state: a map of
  \code{players}, a map of \code{computers}, a map of \code{whiteboards}, and an
  array of \code{chatMessages}.
  \item \code{IPlayer} --- \code{\{name, x, y, anim, readyToConnect,
  videoConnected\}}: everything replicated about an avatar.
  \item \code{IComputer} / \code{IWhiteboard} --- each holds a set of the session
  ids currently connected to it.
  \item The \code{Message} enum --- the remote-procedure opcodes exchanged over
  the socket (\code{UPDATE\_PLAYER}, \code{CONNECT\_TO\_COMPUTER},
  \code{ADD\_CHAT\_MESSAGE}, and so on).
  \item Small enums for \code{ItemType} (\code{CHAIR}, \code{COMPUTER},
  \code{WHITEBOARD}, \code{VENDINGMACHINE}), \code{PlayerBehavior}
  (\code{IDLE}, \code{SITTING}), \code{RoomType} (\code{LOBBY}, \code{PUBLIC},
  \code{CUSTOM}), and day/night background mode.
\end{itemize}

\section{The Authoritative Server}

The server is a Colyseus application. Its entry point creates an Express-backed
Colyseus server and registers three room definitions: the built-in
\code{LobbyRoom} (which tracks the live list of rooms), a persistent
\code{PUBLIC} office room that is never disposed even when empty, and a
\code{CUSTOM} room type for password-protectable private offices whose creation is
broadcast to lobby clients in real time. A Colyseus monitor dashboard is mounted
for live inspection.

Each office is a \code{SkyOffice} room extending \code{Room<OfficeState>}. Its
lifecycle hooks encode the rules of the world:

\begin{description}
  \item[\code{onCreate}] optionally bcrypt-hashes a room password, instantiates
  the state, and populates it with a fixed set of five computers and three
  whiteboards (keyed \code{"0"}\ldots\code{"4"} and \code{"0"}\ldots\code{"2"}),
  then registers all message handlers.
  \item[\code{onAuth}] for a protected room, compares the supplied password
  against the stored hash, rejecting a mismatch with a \code{403}.
  \item[\code{onJoin}] inserts a new \code{Player} keyed by the connection's
  session id and sends that client the room's metadata.
  \item[\code{onLeave}] removes the player and scrubs its session id from every
  computer's and whiteboard's connected-user set, so occupancy stays correct.
\end{description}

State is mutated through the \textbf{command pattern} (\code{@colyseus/command}):
a dispatcher runs small command objects --- \code{PlayerUpdateCommand} (write
\code{x}, \code{y}, \code{anim}), the computer/whiteboard add-user and remove-user
commands, and \code{ChatMessageUpdateCommand}, which appends a chat message and
caps the history at the last 100 so late-joiners see recent context. A few trivial
flag flips (marking a player ready, or relaying a ``stop screen share'' signal)
bypass the dispatcher for directness.

\section{The Client}

The client runs two parallel worlds glued together by a shared Redux store and a
Phaser event emitter: the \textbf{Phaser} game canvas (the animated office) and a
\textbf{React} DOM overlay (dialogs, chat, buttons). Neither engine imports the
other's render tree; they communicate only through the store and the event bus.

\subsection{Scenes}

Phaser organises a game into scenes. The office uses three:

\begin{description}
  \item[Bootstrap] preloads every asset --- the Tiled tile-map, the tilesets, the
  item spritesheets, and the four character spritesheets (\emph{Adam}, \emph{Ash},
  \emph{Lucy}, \emph{Nancy}) --- constructs the network service (opening the
  WebSocket), and pre-checks media permissions.
  \item[Background] draws a decorative parallax backdrop with a sun/moon and
  drifting clouds, switching palette between a day and a night theme.
  \item[Game] is the core scene. It builds the character animations and the
  tile-map with its collision layer, spawns the local avatar and a
  \emph{player-selector} probe, and reads the Tiled object layers to instantiate
  chairs, computers, whiteboards, and vending machines as physics bodies. It then
  registers colliders (avatar versus walls and objects) and two crucial
  \emph{overlap} handlers --- the selector versus interactive items (which drive
  the ``Press E/R'' prompts) and the local avatar versus other avatars (which
  drives proximity calls) --- and wires up the network callbacks.
\end{description}

\subsection{Characters}

Every avatar is a Phaser sprite with a companion container holding a floating name
label and a transient speech bubble; its render depth is derived from its
\code{y}-coordinate so that avatars correctly overlap in the 2.5D view. Three
character classes specialise this base:

\begin{description}
  \item[\code{MyPlayer}] is the locally controlled avatar. Each frame it reads the
  keyboard (WASD or arrows) and the on-screen joystick, moves at a normalised 200
  pixels per second (so diagonals are not faster), plays the matching walk or idle
  animation, and --- crucially --- sends an \code{UPDATE\_PLAYER} message whenever
  its position or animation changes. Pressing \emph{E} near a chair snaps it into a
  seated pose and state; pressing \emph{R} opens a computer, a whiteboard, or a
  vending-machine link depending on what the selector has highlighted.
  \item[\code{OtherPlayer}] is a remote avatar. It does not receive a continuous
  stream of positions; instead it stores a \emph{target} position and interpolates
  toward it each frame at the same 200 px/s, snapping if the gap is large or the
  tab was backgrounded. It also owns the proximity-call trigger.
  \item[\code{PlayerSelector}] is an invisible probe kept one tile ahead of the
  local avatar in its facing direction; whatever interactive item it overlaps
  becomes the current selection that \emph{E}/\emph{R} act upon.
\end{description}

\subsection{The React and Redux overlay}

The React layer is a small state machine driven by Redux: with no room joined it
shows the room-selection dialog; with a room but no login it shows the
name-and-avatar picker; once in the world it renders the chat panel, a prompt to
enable the webcam, and the mobile joystick, overlaying a computer or whiteboard
dialog when the user is at one of those objects. Five Redux slices hold the shared
UI state --- the user/session, the room list, the screen-share session, the
whiteboard URLs, and the chat log --- and it is through these slices that the
Phaser world and the React interface stay in agreement.

\section{The Multiplayer Synchronisation Model}
\label{sec:sky-sync}

The heart of the component is how many moving avatars are kept in sync
efficiently. Colyseus is \emph{authoritative}: it holds the one true
\code{OfficeState}, every client holds a replica, and Colyseus transmits only the
\emph{changed fields} as compact binary diffs --- moving one avatar never
re-sends the whole world. On top of that, the client uses two classic
game-networking techniques (Figure~\ref{fig:sky-sync}):

\begin{description}
  \item[Client-side prediction.] The local avatar moves \emph{immediately} in
  response to input, without waiting for the server. Its new position is then sent
  to the server, so the player never feels input lag.
  \item[Remote interpolation.] Other avatars are not snapped to each incoming
  position; each is eased smoothly from its current position toward the latest
  target, so sparse network updates still produce fluid motion.
\end{description}

\begin{figure}[htbp]
\centering
\begin{tikzpicture}[font=\footnotesize,
  b/.style={draw=voPrimary,fill=voPrimary!8,rounded corners=2pt,minimum height=0.7cm,
            text width=2.3cm,align=center},
  a/.style={-{Stealth[]},thick,voDark},note/.style={font=\scriptsize\itshape,voGray}]
\node[b] (a1) at (0,0) {My client\\(prediction)};
\node[b] (sv) at (4.4,0) {Colyseus\\(authoritative)};
\node[b] (a2) at (9,0) {Other client\\(interpolation)};
\draw[a] (a1) -- node[above,note]{UPDATE\_PLAYER} node[below,note]{\{x,y,anim\}} (sv);
\draw[a] (sv) -- node[above,note]{binary diff} node[below,note]{onChange} (a2);
\draw[a,voGray,densely dashed] (a1.-90) to[out=-90,in=-90] node[below,note]{move now, don't wait} (a1.-90);
\end{tikzpicture}
\caption{Movement synchronisation: local prediction, authoritative relay, remote
interpolation. Only changed fields cross the wire.}
\label{fig:sky-sync}
\end{figure}

The full loop is: \code{MyPlayer} moves locally and sends \code{UPDATE\_PLAYER};
the server's \code{PlayerUpdateCommand} writes the change into that player's
schema; Colyseus diffs it to every \emph{other} client, whose change-callback
updates the corresponding \code{OtherPlayer}'s target; and that \code{OtherPlayer}
interpolates toward it. A single elegant identity choice threads the whole system:
the Colyseus \emph{session id} doubles as the key in every state map \emph{and} as
the PeerJS peer id used for media, so one identifier follows a user everywhere.

\section{Interactive Objects}

The world is populated with objects that blend local responsiveness with
server-authoritative shared state:

\begin{description}
  \item[Chairs] are purely local: pressing \emph{E} snaps the avatar into a
  direction-specific seated pose. Because the resulting position and animation are
  broadcast like any movement, everyone sees you sit --- but the chair itself
  holds no server state.
  \item[Computers] are screen-sharing stations with \emph{server-tracked
  occupancy}. Pressing \emph{R} adds your session id to that computer's connected
  set; the change propagates to all clients, updating an ``N~users'' label, and
  anyone already at that computer receives your shared screen. This makes ``who is
  at this computer'' consistent for everyone.
  \item[Whiteboards] work the same way, and each carries a server-generated random
  id that maps to a hosted collaborative whiteboard shown in an embedded frame, so
  everyone at the same whiteboard shares one board.
  \item[Vending machines] are a decorative touch that open an external link.
\end{description}

The design principle is a clean split: \emph{who is using} a shared object is
authoritative server state (consistent for all), while \emph{what you can interact
with right now} is decided by local physics (responsive, no round-trip).

\section{Proximity Communication}

The office's signature feature is that communication follows position. It is built
on three mechanisms.

\paragraph{Proximity voice and video (PeerJS/WebRTC).} When two avatars physically
overlap, the game's overlap handler attempts a call. To avoid both sides dialling
at once, only the peer with the greater session id initiates, after a short
debounce and only if both peers have enabled their webcam and marked themselves
ready; the other side auto-answers. Encrypted audio and video then flow directly
peer-to-peer. When the avatars separate, the call is torn down --- except inside a
designated private-meeting rectangle, where the call deliberately persists so a
group can gather without staying pixel-perfectly overlapped.

\paragraph{Screen sharing.} A \emph{second}, separate PeerJS peer (its id suffixed
\code{-ss}) isolates screen media from webcam media. When a user shares their
screen at a computer, the manager calls every user currently connected to that
same computer object, and new arrivals are called as their connection propagates
through the shared state; a ``stop sharing'' signal is relayed so remote tiles
disappear cleanly.

\paragraph{Text chat.} A sent message follows two paths from the server at once:
it is appended to the synchronised, capped chat history (so late-joiners see it),
and it is broadcast immediately as a transient speech bubble above the sender's
avatar. The React chat panel shows the scrollback alongside synthetic
``joined''/``left'' system lines.

\section{Integration with the Wider Platform}

The 2D office is a Colyseus-based fork, but it does not stand alone: the backend
is deliberately shaped to drive a spatial client, and a demonstration stack wires
the two together. On the backend side the seams are explicit. The workspace
service exposes an internal \route{/session-config} endpoint that returns exactly
the boot payload a Colyseus room needs --- the workspace, its active desks, and
its active map objects (computers and whiteboards, each already carrying a stable
\code{roomId}) --- and accepts avatar \route{/presence} updates. The room service
consumes avatar position snapshots to compute proximity voice groups
(Section~\ref{sec:be-room}), and both the chat and room services bind their
channels to the same workspace.

These seams are exercised concretely by a one-command \textbf{demonstration
stack} (\code{backend/dev}). A single script brings up the infrastructure
(PostgreSQL, MongoDB, Redis, RabbitMQ) in containers, starts the workspace, chat,
room, and gateway services together with the Colyseus fork and the SkyOffice
client, seeds two workspaces of five members each \emph{consistently across every
service}, mints a demonstration JWT, and prints a URL per user. The workspace
service seeds (through a Flyway migration) a tile map with a meeting-room zone and
open floor, five active desks placed near the spawn so that two users land in the
same proximity voice group, and computers and whiteboards; the chat and room
services seed the matching canonical channel and a shared voice room for the same
members. Opening the two printed URLs in separate browser sessions lets two
avatars walk the same office, meet in proximity to trigger voice, and exchange
chat --- all backed by the real, running microservices rather than by a static
map. Deepening this binding (so that every avatar, desk, and object is served
live from the workspace and every movement feeds the proximity engine in steady
state) is the natural next step, identified in the future-work roadmap of
Chapter~\ref{ch:conclusion}.
